\documentclass[11pt]{article}
\usepackage{amsmath,amssymb,amsthm}
\usepackage[margin=1.1in]{geometry}
\newtheorem{theorem}{Theorem}
\newtheorem{lemma}[theorem]{Lemma}
\newtheorem{corollary}[theorem]{Corollary}
\newtheorem{proposition}[theorem]{Proposition}
\newtheorem{remark}[theorem]{Remark}
\newcommand{\Gm}{\Gamma_{\mathbb C}}
\title{The general Rankin--Selberg pair by polarization}
\author{Samuel Lavery\thanks{Machine-checked companion:
\texttt{RequestProject/RankinSelbergCross.lean}, 13 named declarations, each at axiom
footprint $\{\texttt{propext},\texttt{Classical.choice},\texttt{Quot.sound}\}$ or a subset,
no \texttt{sorry}, no \texttt{axiom}, zero \texttt{sorryAx}.}}
\date{Draft --- \today}

\begin{document}
\maketitle

\begin{abstract}
For an arbitrary pair $f, g$ of level-one cusp forms of one weight $k$, the completed
Rankin--Selberg transform is constructed and its full analytic package proven --- global
self-dual functional equation, polar structure with residue the polarized Petersson
pairing, entire part, Deligne-chart identification, and literal cross Hecke coefficients
--- by \emph{polarization}: the cross object is
$\frac14\bigl[\Lambda_{f+g}-\Lambda_{f-g}+i\,\Lambda_{f+ig}-i\,\Lambda_{f-ig}\bigr]$, and
the compiled single-form engine (the averaged lattice pair with its reflection from the
theta transformation law) applies to each of the four diagonal forms verbatim.  No new
analytic input enters; everything is unconditional and machine-checked.
\end{abstract}

\section{Polarization}

\begin{lemma}[sesquilinear polarization]\label{lem:pol}
For $x,y\in\mathbb C$:
$x\bar y=\tfrac14\bigl[(\|x+y\|^2-\|x-y\|^2)+i(\|x+iy\|^2-\|x-iy\|^2)\bigr]$.
(\texttt{polarization}.)
\end{lemma}

Write $h_1=f+g$, $h_2=f-g$, $h_3=f+ig$, $h_4=f-ig$ --- cusp forms of the same weight
(\texttt{polForm}$_{1\ldots4}$).  Since $q$-expansion is linear
(the compiled expansion-linearity at level one), Lemma~\ref{lem:pol} applied
coefficientwise gives
$a_f(n)\overline{a_g(n)}
=\tfrac14[\,\|a_{h_1}(n)\|^2-\|a_{h_2}(n)\|^2+i\|a_{h_3}(n)\|^2-i\|a_{h_4}(n)\|^2\,]$
(\texttt{cross\_coeff\_polarization}).

\section{The cross pair}

Let $\bar\theta_h$ be the compiled averaged lattice profile of a cusp form $h$ and
$m_h=\|h\|^2$ its Petersson mass.  Define the \emph{cross profile} and \emph{cross mass}
\[
  \Theta_{f,g}=\tfrac14\bigl[(\bar\theta_{h_1}-\bar\theta_{h_2})
    +i(\bar\theta_{h_3}-\bar\theta_{h_4})\bigr],\qquad
  m_{f,g}=\tfrac14\bigl[(m_{h_1}-m_{h_2})+i(m_{h_3}-m_{h_4})\bigr],
\]
the second being the Petersson pairing $\langle f,g\rangle$ in polarized form.  Each
$\bar\theta_{h_j}$ satisfies the weight-one reflection
$\bar\theta_{h_j}(1/t)=t\,\bar\theta_{h_j}(t)$ --- the lattice theta transformation law
carried through the Petersson average, form by form --- so the combination satisfies it
too, and each carries the compiled saddle decay and local integrability.  Hence the data
assemble into a Mathlib weak functional-equation pair with weight $1$, root number $1$,
and both constant terms $m_{f,g}$ (\texttt{crossPair}), every field inherited from the
four compiled diagonal pairs.

\begin{theorem}[the analytic package]\label{thm:pkg}
Let $\Lambda$ be the completed transform of the cross pair.  Then, unconditionally:
$\Lambda(1-s)=\Lambda(s)$ for every $s\in\mathbb C$; $(s-1)\Lambda(s)\to m_{f,g}$ as
$s\to1$; and $\Lambda$ minus its two booked polar terms is entire.
(\texttt{cross\_selfdual\_FE}, \texttt{cross\_residue}, \texttt{cross\_entire}.)
\end{theorem}

\section{Chart and coefficients}

Define the cross Rankin bank as the polarized combination of the four diagonal rank-4
banks, $c_{f,g}=\tfrac14[(c_{h_1}-c_{h_2})+i(c_{h_3}-c_{h_4})]$
(\texttt{crossRankinBank}).

\begin{theorem}[chart identification]\label{thm:chart}
On $2<\Re s$:
$\Lambda(s)=2^{-k}\,\Gm(s)\,\Gm(s+k-1)\,L(c_{f,g},s)$.
\end{theorem}
\begin{proof}
The pair's transform is the Mellin transform of $\Theta_{f,g}-m_{f,g}$; Mellin linearity
(\texttt{mellin\_add$'$}, \texttt{mellin\_sub$'$}, with the four compiled convergences)
splits it into the four diagonal tails; the compiled unpeeled landing evaluates each as
$2^{-k}\Gm(s)\Gm(s+k-1)L(c_{h_j},s)$ --- valid for \emph{arbitrary} cusp forms, no
eigenform input --- and L-series linearity reassembles the combination.
(\texttt{cross\_mellin\_split}, \texttt{LSeries\_crossRankinBank},
\texttt{cross\_lambda\_eq}.)
\end{proof}

\begin{theorem}[coefficient identification]\label{thm:coeff}
$c_{f,g}=\mathbf 1_{\square}\star b_{f,g}$, where
$b_{f,g}(n)=a_f(n)\overline{a_g(n)}\,n^{-(k-1)}$ is the literal cross Hecke datum in
Deligne normalization.
\end{theorem}
\begin{proof}
Each diagonal bank factors as $c_h=\mathbf 1_{\square}\star b_h$ (Möbius inversion
$\zeta\star\mu=\delta$ inside the compiled convolution presentation,
\texttt{rankinBank\_eq\_sqIndicator\_mul}); convolution is bilinear, so the combination
is $\mathbf 1_{\square}\star b_{f,g}$ with $b_{f,g}$ the polarized combination of the
$b_h$ --- which is the stated pointwise datum by coefficient polarization.
(\texttt{crossSquare}, \texttt{crossSquare\_apply}, \texttt{crossRankinBank\_eq\_sq}.)
\end{proof}

\begin{theorem}[rung 6, one theorem]\label{thm:rung6}
For every pair $f,g$ of level-one cusp forms of weight $k\ge0$: the functional equation,
polar structure, entire part, chart identification, and coefficient identification of
Theorems~\ref{thm:pkg}--\ref{thm:coeff} hold simultaneously.
(\texttt{rankin\_selberg\_rung}.)
\end{theorem}

\begin{remark}[register and falsifiers]
Proven and machine-checked: everything above, for arbitrary pairs --- no eigenform
hypothesis anywhere in this note; the diagonal $f=g$ recovers the symmetric-square
companion's pair.  Cited, not re-proved: that the object so constructed is the
automorphic Rankin--Selberg convolution of the classical theory (Rankin, Selberg;
Jacquet for $\mathrm{GL}(2)\times\mathrm{GL}(2)$) --- the identification of coefficients,
chart, functional equation and poles with the classical normalization is exactly what the
theorems state.  A counterexample would be a pair violating $\Lambda(1-s)=\Lambda(s)$, or
a level at which $c_{f,g}$ differs from
$\mathbf 1_\square\star(a_f\bar a_g\,n^{-(k-1)})$ --- each formally stated, so a
counterexample would require a kernel failure.
\end{remark}

\begin{thebibliography}{9}
\bibitem{R} R.~A.~Rankin, \emph{Contributions to the theory of Ramanujan's function
$\tau(n)$ I, II}, Proc.\ Camb.\ Phil.\ Soc.\ \textbf{35} (1939), 351--372.
\bibitem{S} A.~Selberg, \emph{Bemerkungen über eine Dirichletsche Reihe, die mit der
Theorie der Modulformen nahe verbunden ist}, Arch.\ Math.\ Naturvid.\ \textbf{43} (1940),
47--50.
\bibitem{J} H.~Jacquet, \emph{Automorphic forms on $\mathrm{GL}(2)$, Part II}, Lecture
Notes in Math.\ \textbf{278}, Springer (1972).
\end{thebibliography}

\end{document}
